% ch-timeintegration.tex -- numerical method: theta method, Newton, coupling and the adaptive step
\chapter{Time integration}
\label{ch:timeintegration}

Every equation Yuzen solves is cast into the same shape, advanced with the same
$\theta$ method, and --- where it is non-linear --- driven to the solution of
the implicit system by the same outer Newton iteration. This chapter describes
that machinery: the form, the iteration, the Jacobian, the way the thickness
and the concentration are coupled, and the adaptive time step. The keys
mentioned here are documented in \cref{sec:ref-numerics:height,%
sec:ref-numerics:tsc}; what happens when the iteration will not converge is the
subject of \cref{ch:convergence}.

\section{The form every equation is cast into}
\label{sec:time:form}

Every equation is rearranged into

\[
  \frac{\partial\varphi}{\partial t} = -A\,\varphi - B
\]

before it is handed to the solver (\file{numerics/ImplicitSolver.h}). $A$ is
the part proportional to the unknown, and is a matrix; $B$ is everything else,
and is a vector. Time is advanced with the $\theta$ method, where
$\theta = 0$, $0.5$ and $1$ are forward Euler, Crank--Nicolson and backward
Euler respectively. \textbf{Both the thickness equation and the concentration
equation are built with Crank--Nicolson}, $\theta = 0.5$.

The system actually solved for the correction $\delta\varphi$ is

\begin{equation}\label{eq:nk-update}
  (I + \theta\,\Delta t\,J)\,\delta\varphi
    = \varphi_{\mathrm{old}} - \varphi_{\mathrm{guess}}
      - \Delta t\left[\theta\,A\cdot\varphi_{\mathrm{guess}}
        + (1-\theta)\,A\cdot\varphi_{\mathrm{old}} + B\right],
\end{equation}

with $J$ the Jacobian of \cref{sec:time:jacobian}. The right-hand side of
\cref{eq:nk-update} is the residual, and it is worth noting now that it is
also the origin of the correction floor of \cref{sec:convergence:floor}.

\begin{warn}[A term in $B$ has to be averaged in time as well]
  The $\theta$ weighting of \cref{eq:nk-update} applies to $A$. A term put into
  $B$ gets it only if $B$ itself is formed as the same average, and the
  well-mixed thickness equation does so: $B = (1-\theta)\,B(\varphi_{\mathrm{old}})
  + \theta\,B(\varphi_{\mathrm{guess}})$ for the evaporation, and the liquid
  supply of \key{Physics.Source} is taken at the middle of the step
  (\code{SetMatrixA} and \code{BeforeNKLoop} of
  \file{HeightEquationBase.h}). Taken at the start of the step alone, $B$
  makes the scheme first order in $\Delta t$ for that term. When you add a new term, put it in $A$
  if it can be written as a coefficient times the unknown --- the disjoining
  pressure $\Pihat(H)$ written as $(\Pihat/H)\cdot H$ is the example --- or
  give it the same average in $B$.
\end{warn}

\subsection{How the product with the matrix is formed}
\label{sec:time:apply}

The two products $A\cdot\varphi_{\mathrm{guess}}$ and
$A\cdot\varphi_{\mathrm{old}}$ of \cref{eq:nk-update} are, for the well-mixed
thickness equation, not all formed as the product of the entries of $A$ with
$\varphi$ (\code{ApplyA} of \file{HeightEquationBase.h}). The capillary
pressure and the vertical gravity form theirs face by face
(\code{AddConservative} of their kernels): the flux through a face is its
mobility --- the one $A$ was assembled with --- times a difference of the film
(the third difference across the face, and its cross term in 2-D, for the
capillary pressure; $h_b - h_a$ and the substrate term for the gravity),
written as differences of neighboring values, and the two cells that share
the face take it from the same values in the same order, so one loses exactly
what the other gains. In exact arithmetic this is the product with $A$; in
floating point the sum over the cells is the flux through the edges of the
domain to rounding of the fluxes. The product with the entries is not: the
entries of a fourth-order term are of order $D/\Delta x^{4}$, and the sum of
the products over the cells would come out as
$\epsmach\,\lVert A\rVert\,\lvert\varphi\rvert$ per cell --- on a film
\SI{1.2}{\milli\metre} thick over cells \SI{3.3}{\micro\metre} wide, as in
\file{nozzle\_1d\_rotating}, a volume error of $2\times10^{-7}$ of the film
over the run (\cref{sec:verification:analytic}). The other terms --- the in-plane body
forces, the Marangoni stress, the disjoining pressure and the sideways
outflow of the 1-D centrifugal term --- keep the product with their entries,
which are of order $1/\Delta x$ or $1/\Delta x^{2}$; so do the rows of the
boundary conditions. $A$ itself, which the Jacobian starts from, still holds
every term.

The flux of an edge face is zero only if the ghosts make it so. After every
linear solve the rows of the boundary conditions are satisfied exactly by the
correction (\code{CRS\_Matrix::SatisfyBCRows}), not merely to the tolerance of
the solve, and a \val{"Symmetry"} ghost is then the exact mirror of the layer
inside it, and a \val{"Periodic"} one the exact copy of the other end: the
third difference across a mirrored edge face is exactly zero, where a ghost
off its mirror by the tolerance of the solve would leak as much as the
product above.

\section{The Newton--Kantorovich outer iteration}
\label{sec:time:newton}

The thickness equation is non-linear because its mobility $h^{3}/(3\mu(c))$
depends on the unknown. Yuzen therefore iterates within a single time step
until the non-linearity itself is solved, rather than lagging it
(\file{numerics/NewtonKantorovich.cpp}). One iteration has five stages:

\begin{enumerate}
  \item \code{EveryNKLoop()} --- update the fields from the current guess
    $\varphi_{\mathrm{guess}}$ and solve the coupled side, the concentration.
  \item \code{SetMatrixA()} --- assemble the matrix $A$.
  \item \code{SetMatrixJacobi()} --- assemble the Jacobian, once every
    \key{JacobianUpdateInterval} iterations; then \code{MakeMatrix\_IA()} forms
    $I + \theta\,\Delta t\,J$.
  \item the linear solve (\cref{ch:linearsolver}), which yields the correction
    $\delta\varphi$.
  \item $\varphi_{\mathrm{guess}} \mathrel{+}= \alpha\,\delta\varphi$, with
    $\alpha = 0.5$ only when the correction has grown since the previous
    iteration, and $\alpha = 1$ otherwise.
\end{enumerate}

The iteration starts from the film extrapolated linearly from the last two
steps, $\varphi^n + (\Delta t/\Delta t_{\mathrm{old}})(\varphi^n -
\varphi^{n-1})$, unless that would leave a cell that has liquid without any;
an attempt from it that fails is made once more from $\varphi^n$, at the same
step, before the step is cut or the run stopped (a source that turns on, a
nozzle, is not in the last two steps). The linear solve of each pass is
solved only as far as the progress of the iteration needs
(\cref{sec:linsolve:ew}). The two cut the Newton iterations by up to a
third and the linear iterations by up to 38\,\% on the cases measured, and
leave the converged answer where it was: the convergence tests are relative
to the first correction of the step, whatever the start
(\cref{sec:convergence:tests}).

The concentration equation, and the thickness equation of the \val{"CZ"} model,
are weakly non-linear. They freeze their coefficients and solve a linear
system --- the semi-implicit path of \file{numerics/SemiImplicit.cpp} ---
with no outer iteration. Frozen at the start of the step, that is first order
in $\Delta t$ whatever $\theta$ is, so the \val{"CZ"} model solves each step
twice: once with the coefficients of the start of the step, which predicts the
end state, and again from the same start with every coefficient and every term
of $B$ formed at the midpoint between the two. That is the implicit midpoint
rule, second order. The thickness equation does this with the concentration
extrapolated to the midpoint from the last two steps. The velocity the solute
is carried with is then formed from the film halfway to the end state, not to
the prediction, so that it carries the flux the film moved with. Without the
Marangoni term, the concentration equation is solved once after that, with a
predicted concentration of its own. With it, the film and the concentration
are solved again, one stage each, at the midpoint concentration the last pass
found, with Aitken relaxation, until that midpoint moves by less than a tenth
of the extrapolation's miss (at most eight passes; \file{JobInfo.txt} reports
how many were taken). That is the implicit midpoint rule for the coupled
pair. On a flat film the gain of a plain pass exceeds 1 only where a wave of
the model grows faster than $2/\Delta t$ (\cref{sec:verification:czpasses}
derives that gain and measures it, and the passes the runs take); a step whose passes do not settle within eight has failed,
as one whose linear solve does not converge: the time-step control cuts it,
and at a fixed step the run stops. The
well-mixed concentration equation is solved inside the Newton loop, on the
current guess of the end state, so its Crank--Nicolson average is already that
of the two ends; once the loop has converged it is solved one last time, with
the face fluxes of the step the film took (\cref{sec:time:coupling}).

\section{The Jacobian}
\label{sec:time:jacobian}

Each physical term contributes its own \emph{analytic} Jacobian through
\code{SetJacobiMatrix()}; this is part of the contract every term implements
(\file{PhysicsTermBase.h}). The difference coefficients are separated from the
field-dependent factors: \code{SetLinearCoef()} builds the coefficients that
depend only on the mesh and the scheme, once, while the part that depends on
the fields lives in \code{NonlinearTerm()} and its derivative
\code{NonlinearJacobi()}.

\textbf{The Jacobian is allowed to be approximate.} An approximate Jacobian
changes how fast the iteration converges; it does not change where it converges
to. What fixes the converged answer is the residual \code{Rhs} of
\cref{eq:nk-update}, not the matrix multiplying the correction. Within the
thickness equation Yuzen exploits this in one place only, the modified Newton
method below; the terms themselves do not need it.

\begin{itemize}
  \item The Jacobian starts as a copy of $A$ and each term adds the
    derivative of its coefficients. For the gravity and centrifugal terms the
    gradient side is linear in $h$, so $A$ already holds its derivative and
    the term adds only that of the mobility; the result is exact, and
    \file{test\_jacobian.cpp} checks every term against a finite difference
    (\cref{sec:verification:jacobian}).
  \item \key{Numerics.HeightEquation.JacobianUpdateInterval} greater than 1 turns
    the scheme into a modified Newton method: the Jacobian is rebuilt only once
    every few iterations, and \code{\_crsIA} can be held with it, since it
    depends on nothing but \code{\_crsJacobi} and $\Delta t$. This is worth
    doing because the Jacobian assembly is the single largest item in an
    iteration. The 1-D breakdown is \code{SetMatrixA} \SI{19}{\percent},
    \code{SetMatrixJacobi} \SI{43}{\percent}, \code{IA + RHS}
    \SI{2}{\percent} and the linear solve \SI{37}{\percent}.
\end{itemize}

\textbf{A gradient side that is not linear needs its own derivative.} The
disjoining pressure goes like $\Pi \sim h^{-n}$ and is extremely stiff, and it
is written into $A$ as $(\hat\Pi/h)\,h$; its term therefore adds the
derivative of the mobility and $\hat\Pi' - \hat\Pi/h$ on the gradient side
(\cref{sec:terms:disjoining}). Adding all of $\hat\Pi'$ would count
$\hat\Pi/h$ twice and leave the Jacobian \SI{50}{\percent} off; with the
correct derivative the disjoining-pressure cases take 30--55 \% fewer Newton
iterations.

\section{Coupling the thickness and the concentration}
\label{sec:time:coupling}

Thickness and concentration are solved as a block Gauss--Seidel iteration
inside the outer Newton loop.

\begin{plain}
within a time step, for k = 0, 1, 2, ...
    form the face volume fluxes F of h^k
    solve the concentration equation over dt   -> c^{k+1}
        (fluxes and coefficients frozen at h^k, the film at the end h^k)
    cap the c of the coefficients at 1         (ClampConc)
    take one Newton step on the thickness      -> h^{k+1}
        (Jacobian is dF_h/dh only)
once converged, to h^1:
    the fluxes F at h^1, made those of the step h^0 -> h^1
    solve the concentration equation once more -> c^1
    bring c^1 into [0, 1] without cutting solute
\end{plain}

The off-diagonal blocks $\partial F_h/\partial c$ and $\partial F_c/\partial h$
are dropped, for exactly the reason given in \cref{sec:time:jacobian}: they
slow the convergence down but \textbf{do not move the converged answer}. If the
loop converges to $\delta h \to 0$ together with $c^{k+1} = c^{k}$, then the
thickness equation and the concentration equation are satisfied
\emph{simultaneously}, and that state is precisely the fully implicit coupled
discretization of the two.

The last solve completes the step. Within the loop each pass carries the
solute with the fluxes of its own film, which give that film only once the
iteration has converged; at the end, the fluxes at the converged film are
made, cell by cell, those of the step the film took (the residual of the
step there is put into the inner faces), and the concentration is solved with
them once more, so that the solute of every cell changes by what crosses its
faces and what the sources bring, and a uniform concentration stays uniform
(\cref{sec:disc:wm-conservative}). The film is not changed by it. The
concentration it gives differs from the last pass's by the tolerance the
loop converged to.

That argument is also what dictates the convergence test. It requires both
conditions, so the test must be made on \textbf{the larger of the two
corrections},

\[
  \max\bigl(\,\lVert\delta h\rVert_{\infty},\
             \lVert\Delta c\rVert_{\infty}\,\bigr).
\]

\begin{warn}[Do not test on the thickness correction alone]
  $\lVert\delta h\rVert$ can become small while the concentration is still
  moving, so a test that leaves the concentration out lets the iteration exit
  early. What it exits to is a state that satisfies the thickness equation
  only, not the fully implicit coupled solution. Propping the iteration up with
  a minimum iteration count instead is not a fix: it decides the accuracy with
  a number that has nothing to do with the problem.
\end{warn}

The concentration correction is measured \textbf{after} the clamp
(\code{ClampConc}), because the clamped value is the one that enters the
coefficients of the physical terms, and therefore the one that decides whether
the coupling has settled. A model that does not report a coupling leaves the
reported value at zero --- the \val{"CZ"} model solves the thickness equation
semi-implicitly too and has no outer iteration at all --- so there the test
reduces to the thickness. \key{NewtonMinIterations} can force a minimum iteration count,
but since the concentration correction is already part of the test, the default
of \val{1} is normally the right value.

To see whether it is the thickness or the coupling that is holding convergence
up, set \env{YUZEN\_DEBUG\_RESIDUAL}. Each outer iteration then prints a line
to standard error in the form
\code{[resid] <iteration> <thickness correction> <concentration correction>}.

\section{The adaptive time step}
\label{sec:time:dt}

The step changes only when \key{Numerics.TimeStepControl.Enable} is
\val{true}. It is decided by the thickness equation; the concentration equation
follows whatever the thickness equation asks for.

\subsection{The local truncation error}

The routine \code{LocalTruncationCriterion} of \code{C\_ImplicitSolver} builds
the estimate from two consecutive time differences, weighing the change over
the current step against the change over the previous one:

\begin{equation}\label{eq:lte}
  \mathrm{Ek}_i = 2\,\frac{\Delta t}{\Delta t_0}\,
    \frac{\Delta t_0\,h^{+}_{i} + \Delta t\,h^{-}_{i}
          - (\Delta t_0+\Delta t)\,h^{0}_{i}}
         {(\Delta t_0+\Delta t)\,w_i},
  \qquad
  \mathrm{criterion} = 1 - \frac{\max\lvert\mathrm{Ek}\rvert}
                                {\mathrm{Criterion}}.
\end{equation}

Here $\Delta t_0$ is the step last taken and $\Delta t$ the current one; for
the first step, $\Delta t_0$ is \key{InitialDt}, so the first step is judged
too. A positive \code{criterion} in \cref{eq:lte} accepts the step; a
negative one shrinks it. A step whose solve fails --- the Newton iteration
of the well-mixed film, not converged in \key{NewtonMaxIterations} or
stagnated far from convergence (\cref{sec:convergence:stagnation}), or a
linear solve of the film or of the
concentration in either model --- is shrunk in the same way and solved
again; with the step control off it stops the run
(\cref{sec:ref-numerics:tsc}). $\Delta t_0$ must be the step that produced $h^{-}$:
taking instead the value of the step before its last change, or the rejected
step on a retry, leaves a first-order term in the estimate and holds the step
down --- in 1-D it does not grow again after a cut. \key{Criterion} is
therefore not the error itself: a leveling film at \key{Criterion}
$= 10^{-5}$ takes 69 steps for an error in its decay of $2.9\times10^{-4}$
(the estimate with the wrong $\Delta t_0$ takes 385 for $1.1\times10^{-5}$).
Choose \key{Criterion} by the error you need.

The weight $w_i$ that normalizes the error is

\begin{equation}\label{eq:lte-weight}
  w_i = \max\bigl(\,\lvert h^{0}_{i}\rvert,\
                   \mathrm{ScaleFloor}\times\max\lvert h^{0}\rvert\,\bigr),
\end{equation}

so \key{ScaleFloor} chooses what ``error'' means. \val{0}, the default, divides
by the local thickness and gives a completely local relative error; \val{1}
normalizes by the global maximum alone and gives something close to an absolute
error; values in between judge the body of the film relatively while refusing
to over-weight the thin regions. Read as a sentence, \cref{eq:lte-weight} says
``anywhere thinner than \key{ScaleFloor} times the maximum thickness is not
judged more harshly than that'', which also serves as the division-by-zero guard
in a dry region, $h \to 0$.

The default has a consequence worth knowing about. A precursor film at
$\Hstar = 0.02$ is a fiftieth of the body of a drop at $H \sim 1$, so with
\key{ScaleFloor} left at \val{0} a wobble inside the precursor film weighs
fifty times as much as the same wobble in the body. \textbf{This is still not
the usual reason a contact-line calculation sees its step collapse}; suspect
the convergence test first (\cref{sec:convergence:absolute}).

\subsection{Growing and shrinking}

\begin{yzntable}{How the step is changed\label{tab:dt-change}}{@{}l L{0.34\textwidth} L{0.30\textwidth}@{}}
\textsf{\bfseries Direction} & \textsf{\bfseries Condition} &
\textsf{\bfseries Action} \\ \midrule
shrink & \code{criterion <= 0} & \code{dt *= ChangeFactor} \\
grow   & \code{criterion > 0.9} for \key{KeepSteps} consecutive steps &
         \code{dt /= ChangeFactor} \\
\end{yzntable}

A shrink acts at once, on the step being retried; a lengthened step applies
from the next step, and the clock advances by the step that was solved.

The step is then clipped from both sides. The upper bound is the smaller of
\key{MaxDtFactor} $\times$ (smallest mesh width)$^{4}$ and the smallest output
interval; the lower bound is \key{MinDt}.

The output interval appears in the upper bound because a step that straddles an
output time would shift the time at which the output is actually written.

The lower bound behaves differently: \textbf{the moment a step smaller than
\key{MinDt} is requested the run aborts}, rather than continuing with the
floor. This is deliberate. A calculation that has shrunk that far is not
suffering from a step-size problem, so lowering \key{MinDt} and letting it grind
on achieves nothing. In practice a step size collapsing in a contact-line case
is the convergence test, not the truncation-error estimator ---
\cref{sec:convergence:absolute} explains the mechanism, and it is the first
thing to check.
